\chapter{Probabilidades++}

\formal
\section{Axiomas de probabilidad}

En el capítulo anterior vimos que la probabilidad se define como una función que cumple ciertas condiciones. Definámosla bien de una vez por todas.
\begin{definition}[Probabilidad]
\labdef{probabilidad_2}
Sea \(S\) un espacio muestral y \(E\) un evento cualquiera de éste. Se llamará \emph{función de probabilidad} sobre el espacio muestral \(S\) a \(\prob{(E)}\) si satisface los siguientes axiomas:
\begin{enumerate}
    \item \(\prob{(E)} \geq 0\)
    \item \(\prob{(S)} = 1\)
    \item Si para los eventos \(E_1, E_2, \ldots\), \(E_i \cap E_j = \emptyset\) para \(i \neq j\), entonces \(\prob{\left(\bigcup_{i} E_i\right)} = \sum_{i} \prob{(E_i)}\).
\end{enumerate}
\end{definition}

De esta definición se deducen los teoremas que hemos estado usando en la sección anterior. Por ejemplo, 
\begin{theorem}
    \(\prob{(\emptyset)} = 0\).
\end{theorem}

\begin{proof}
    \(S \cup \emptyset = S\) y \(S \cap \emptyset = \emptyset\). Por el axioma 3, \[\prob{(S\cup \emptyset)} = \prob{(S)} + \prob{(\emptyset)},\] pero por el axioma 2, \(\prob{(S)} = 1\). Por lo tanto, \[\prob{(\emptyset)} = 0.\]
\end{proof}

\informal

Podemos demostrar los teoremas de probabilidad usando los axiomas de manera similar. Como dicen los libros, queda como ejercicio al lector.

\section{Teorema de Bayes}
\labsec{teo_bayes}

Nos han descrito a Esteban de la siguiente manera:
\begin{itemize}
\item Es organizado.
\item Le gusta el orden.
\item Es callado, pero afable.
\item Es servicial.
\end{itemize}

La pregunta es si Esteban será un bibliotecario o un campesino (no hay más opciones). En una de las secciones, aproximadamente el \qty{85}{\percent} pensó que Esteban sería un bibliotecario versus un \qty{15}{\percent} que pensó que sería un campesino.

Todo eso está muy bien, pero hay una información que no hemos considerado y que, francamente, debiéramos haber tenido en cuenta: hay muchos más campesinos que bibliotecarios. Como ejemplo, podemos asumir que la razón entre bibliotecarios y campesinos es \(1:20\) (en realidad debe ser algo como \(1:100\)).

Siendo así, podemos tomar una muestra\sidenote{más adelante en el curso definiremos qué es eso} que siga esa proporción (\reffig{muestra_bibliotecarios_campesinos}). 

\begin{figure*}
    \begin{tikzpicture}[scale=0.5]
    \foreach \i in {1, ..., 10}{
        \foreach \j in {1, ..., 20}{
            \node[text=Peach] at (\j, \i) {\faUser};
            \node[text=RoyalBlue] at (0, \i) {\faUser};
        }
    }
    \end{tikzpicture}
    \caption{Muestra de bibliotecarios y campesinos.}
    \labfig{muestra_bibliotecarios_campesinos}
\end{figure*}

Supongamos que aproximadamente un \qty{40}{\percent} de los bibliotecaros siguen la descripción de Esteban, pero sólo un \qty{10}{\percent} de los campesinos la siguen. Si fuera así, de nuestros 10 bibliotecarios, 4 serían como Esteban y de los 200 campesinos, 20 serían como Esteban (\reffig{bibl_camp_esteban}).

\begin{figure*}
    \begin{tikzpicture}[scale=0.5]
    \foreach \i in {1, ..., 10}{
        \foreach \j in {1, ..., 20}{
            \node[text=Peach] at (\j, \i) {\faUser};
            \node[text=RoyalBlue] at (0, \i) {\faUser};
        }
    }
    \fill[white, opacity=0.6] (-0.5,0) rectangle (0.5,6.5);
    \fill[white, opacity=0.6] (0.5,0) rectangle (20.5,9.5);
    \end{tikzpicture}
    \caption{Bibliotecarios y campesinos como Esteban.}
    \labfig{bibl_camp_esteban}
\end{figure*}